\section{Introduction}
\label{sec:introduction}

Sparse linear regression---finding the few predictors that actually matter in a potentially large feature space---comes up everywhere: genomics \citep{tibshirani1996regression}, signal processing, economics, neuroscience.
In many of these fields, knowing \emph{which} variables are active matters as much as getting the prediction right.

\paragraph{Two camps, one problem.}
Practitioners typically reach for one of two toolboxes.
Penalized methods---Lasso \citep{tibshirani1996regression}, Ridge \citep{hoerl1970ridge}, Elastic Net \citep{zou2005regularization}---are fast and well-understood, but they output point estimates only.
If you need error bars on your coefficients, you are out of luck.
Bayesian priors like the Horseshoe \citep{carvalho2010horseshoe} and Spike-and-Slab \citep{mitchell1988bayesian} give you full posteriors, credible intervals, and a principled inclusion probability---at the cost of MCMC sampling that can be orders of magnitude slower.

\paragraph{What is missing.}
Both camps have been studied extensively in isolation.
But when a practitioner asks ``should I use Lasso or Horseshoe for \emph{my} problem?'', the existing literature is surprisingly unhelpful.
Most comparisons fix the correlation structure (usually independent features), test at one or two SNR levels, and skip the calibration question entirely.
Real data is rarely that clean: features are correlated, signals can be weak, and dimensions vary.

\paragraph{What we do.}
We run a large-scale benchmark that varies \emph{all three axes simultaneously}:
\begin{enumerate}[nosep]
    \item \textbf{Correlation}: independent, block-diagonal, and Toeplitz designs with $\rho$ from 0 to 0.9;
    \item \textbf{Signal strength}: SNR from 0.5 (buried in noise) to 5 (clearly visible);
    \item \textbf{Dimensionality}: $p \in \{20, 50, 100\}$.
\end{enumerate}
We test six methods across this grid with five seeds each, evaluating prediction (MSE), estimation ($L_2$ error), selection ($F_1$), and---for Bayesian methods---posterior calibration (coverage and interval width).
The result is over 2{,}600 experiments, all reproducible from a single codebase.\footnote{\url{https://github.com/xiao98/sparse-bayesian-regression-bench}}

\paragraph{Main takeaways.}
\begin{itemize}[nosep]
    \item Bayesian methods predict better overall, and the gap widens under high correlation.
    \item Horseshoe intervals are well-calibrated; Spike-and-Slab's continuous relaxation leads to slight under-coverage.
    \item When you do not need posteriors, Lasso remains hard to beat for speed and selection quality.
    \item High correlation ($\rho > 0.6$) breaks Lasso's variable selection; Elastic Net and Bayesian methods are more robust.
\end{itemize}

\paragraph{Outline.}
Section~\ref{sec:related_work} covers related work.
Sections~\ref{sec:methods}--\ref{sec:experimental_setup} describe the methods and setup.
Section~\ref{sec:results} presents results, and Section~\ref{sec:discussion} discusses practical implications.
